\documentclass[a4paper,10pt]{article}
\usepackage[margin=17mm]{geometry}
\usepackage{fontspec,amsmath,array,tabularx,xcolor}
\setmainfont{Noto Sans CJK JP}
\setsansfont{Noto Sans CJK JP}
\XeTeXlinebreaklocale "ja"
\XeTeXlinebreakskip=0pt plus 1pt
\definecolor{navy}{HTML}{183248}
\definecolor{blue}{HTML}{276698}
\definecolor{green}{HTML}{24734A}
\definecolor{orange}{HTML}{B45B21}
\setlength{\parindent}{0pt}
\setlength{\parskip}{5pt}
\renewcommand{\arraystretch}{1.34}
\newcommand{\ttl}[2]{{\color{navy}\Large\bfseries #1}\par{\small #2}\par\vspace{2pt}\hrule\vspace{7pt}}
\newcommand{\sub}[1]{\par{\color{navy}\large\bfseries #1}\par}
\newcommand{\boxwork}[1]{\par\noindent\fbox{\parbox[t][#1][t]{0.965\linewidth}{\vspace{1pt}\hspace{2mm}{\color{gray}\small 途中式・検算}}}\par}
\newcommand{\note}[1]{{\small\color{blue}#1}\par}
\begin{document}
\ttl{第05回 トラバース測量 I｜教員用解答（改3）}{採点10点／同じ内部数値からPDF・Excel・JSONを再生成}
\sub{問1 (1)　角閉合差と等配分（2点）}
観測和は540°01′00″、理論和は$(5-2)\times180°=540°$。
\[
\varepsilon_\beta=\sum\beta_{\rm obs}-(n-2)180°=+60″,
\qquad v_\beta=-\frac{\varepsilon_\beta}{n}=-12″.
\]
授業用許容±80″に対し $|+60″|\le80″$ なので許容内。各角等精度という条件で等配分する。許容外なら機械的に配分せず、観測・記録を点検する。
\begin{center}\small\begin{tabular}{|l|l|l|l|}
\hline
\textbf{点} & \textbf{観測β} & \textbf{補正} & \textbf{調整β} \\ \hline
A & 100°12′12″ & −12″ & 100°12′00″ \\ \hline
B & 118°35′52″ & −12″ & 118°35′40″ \\ \hline
C & 95°48′32″ & −12″ & 95°48′20″ \\ \hline
D & 128°10′42″ & −12″ & 128°10′30″ \\ \hline
E & 97°13′42″ & −12″ & 97°13′30″ \\ \hline
和 & 540°01′00″ & −60″ & 540°00′00″ \\ \hline
\end{tabular}\end{center}
\note{採点：観測和・εβと許容判定0.5点、補正−12″ 0.5点、調整角と総和1.0点。}
\sub{問1 (2)　方向角（2.5点）}
後視の方向角は前のT＋180°。そこからβを時計回りに加え、0°以上360°未満に整える。
\[T_{\rm next}=\operatorname{mod}(T_{\rm prev}+180°+\beta,\ 360°)\]
\begin{center}\small\begin{tabular}{|l|l|l|}
\hline
\textbf{測線} & \textbf{後視方向角＋調整β（整える前）} & \textbf{方向角T} \\ \hline
AB & 既知 & 115°20′30″ \\ \hline
BC & 413°56′10″ & 53°56′10″ \\ \hline
CD & 329°44′30″ & 329°44′30″ \\ \hline
DE & 637°55′00″ & 277°55′00″ \\ \hline
EA & 555°08′30″ & 195°08′30″ \\ \hline
AB 検算 & 475°20′30″ & 115°20′30″ \\ \hline
\end{tabular}\end{center}
\note{採点：BC・CD・DE・EA各0.5点（±2″）、ABに戻る検算0.5点。}
{\small 旧版は5点の数値を維持し、角許容80″を授業条件として明記した。講義4点の許容40″を流用すると60″は許容外になる。両者はいずれも規程値ではない。}
\newpage
\ttl{問1 (3)　成分計算と総和（1.5点）}{X＝北・Y＝東／計算途中では丸めない}
\[\Delta X=S\cos T,\qquad \Delta Y=S\sin T\]
\begin{center}\small\begin{tabular}{|l|l|l|l|l|}
\hline
\textbf{測線} & \textbf{S (m)} & \textbf{方向角T} & \textbf{ΔX (m)} & \textbf{ΔY (m)} \\ \hline
AB & 85.000 & 115°20′30″ & −36.381 & +76.821 \\ \hline
BC & 110.000 & 53°56′10″ & +64.756 & +88.920 \\ \hline
CD & 96.000 & 329°44′30″ & +82.921 & −48.374 \\ \hline
DE & 84.894 & 277°55′00″ & +11.693 & −84.085 \\ \hline
EA & 127.412 & 195°08′30″ & −122.989 & −33.281 \\ \hline
内部総和 & 503.306 &  & −0.001 & +0.000 \\ \hline
\end{tabular}\end{center}
BCの計算例：$110.000\cos(53°56′10″)=+64.755569\ldots$ m、
$110.000\sin(53°56′10″)=+88.919717\ldots$ m。

\textbf{総和の表示に注意}

未丸めの内部値から計算すると、
ΣΔX＝-0.000501822105 m、ΣΔY＝+0.000172928216 m。

したがって表の内部総和は0.001 m表示で −0.001 m、0.000 mとなる。各辺を先に0.001 mへ丸めて足すと0.000 m、＋0.001 mとなり、違う値になる。検算・補正には内部総和を使う。

\note{採点：BC・CD・EAの各辺0.5点、計1.5点。各成分の符号が正しく、±0.002 mを目安に評価する。総和は検算課題であり、得点内訳に含めない（本演習は計10点）。表示値を足した合計が出た場合は、内部値総和との違いを説明する。}
\textbf{小さい閉合差の解釈}

距離は合成教材用に設定してある。小さい座標閉合差や大きい閉合比だけを、実測の精度が高い証拠にしてはいけない。角の点検、方向角の一周検算、象限と符号をそれぞれ確認する。閉合比とBowditch補正は次回の内容で、本演習の採点対象外。

\textbf{授業で拾う誤り}

・DEG設定のまま度で計算する。ExcelはCOS(RADIANS(T))を使う。

・BCは北東で両成分が正、CDは北西でΔYが負。

・観測角で方向角を作らず、調整角を用いる。

・AB・DEの表示成分は確認用。総和を求めるときはそれらも元の距離と方向角から再計算する。

\sub{発展（加点なし）}
観測角のままで一周すると角和が理論値より60″大きいので、ABの方向角は115°21′30″となり、既知値から＋60″ずれる。
\newpage
\ttl{問2　座標の逆計算（4点）}{距離・方向角を求め、順計算で検算する}
\sub{(1) 座標差と象限（1点）}
\[\Delta X=1580-1500=+80.000\ \mathrm{m},\qquad
\Delta Y=2150-2300=-150.000\ \mathrm{m}\]
ΔX＞0、ΔY＜0なので、講義の象限表記では第4象限（北西）。
\note{採点：符号付き座標差0.5点、象限・向き0.5点。}
\sub{(2) 距離（1点）}
\[S_{PQ}=\sqrt{80^2+(-150)^2}=\sqrt{28900}=170.000\ \mathrm{m}\]
\note{採点：式と値1.0点。8：15：17の直角三角形。}
\sub{(3) 方向角（1.5点）}
\[\arctan(\Delta Y/\Delta X)=\arctan(-1.875)=-61°55′39″\]
第4象限なので360°を加え、$T_{PQ}=298°04′21″$。

内部値は298.07248693585296°。1秒表示へ丸める前の角で順計算すると、ΔX＝＋80.000 m、ΔY＝−150.000 m、Q＝(1580.000, 2150.000) mへ戻る。1秒表示角を再入力すると小さな丸め差が出るため、元の内部値と混ぜない。
\note{採点：arctanの値0.5点、＋360° 0.5点、Tの値（±2″）0.5点。}
\sub{(4) 反方向角（0.5点）}
\[T_{QP}=\operatorname{mod}(T_{PQ}+180°,360°)=118°04′21″\]
\note{採点：0.5点。}
\sub{逆計算の境界条件}
arctan(ΔY/ΔX)だけでは象限を区別できず、ΔX＝0で割れない。プログラムなら象限対応のatan2を使う。Pythonではatan2(ΔY,ΔX)、ExcelではATAN2(ΔX,ΔY)。同一点は距離0、方向角は未定義とする。

\vfill
{\small 再生成原稿：notes/session05\_20261117/build\_numeric.py。数値例は合成教材。距離は換算済み平面計算用距離。X北・Y東、Tは北から時計回り。}
\end{document}
